\newpage
\section{第16章：子类型化的元理论}

\begin{taplauthorsolutionentrybox}
\phantomsection
\label{sol:author:16.1.2}
\textbf{16.1.2 解答}

\emph{第一项。}对 \(S\) 的结构作归纳。若 \(S=\tapltype{Top}\)，直接使用
\textsc{S-Top}；若 \(S=S_1\to S_2\)，先根据归纳假设分别构造
\(S_1\subtype S_1\) 与 \(S_2\subtype S_2\) 的无自反规则推导，再应用
\textsc{S-Arrow}。记录类型的情形相同：先对各字段类型使用归纳假设，再应用
\textsc{S-Rcd}。因此，任何 \(S\subtype S\) 的推导都不必使用
\textsc{S-Refl}。

\emph{第二项。}先利用第一项，把原推导中的每次 \textsc{S-Refl} 都换成不使用
自反规则的推导。接下来只需证明：任何不含 \textsc{S-Refl} 的
\(S\subtype T\) 推导，都能进一步改造成不含 \textsc{S-Trans} 的推导。

把待改造的推导视为一棵树，并对它的大小作强归纳。这里按大小而不按推导结构
归纳，是因为在箭头和记录情形中，我们会把原推导的若干部分重新组合成新的推导；
这些新推导不是原推导中现成的子树，但只要规模严格小于原推导，强归纳假设仍然
适用。

若原推导的末条规则不是 \textsc{S-Trans}，就分别对它的每个直接子推导使用
归纳假设，消去其中的传递规则。末条规则本身也不是传递规则，因此重新组合之后，
整个推导便不再含有 \textsc{S-Trans}。

下面只需处理末条规则为 \textsc{S-Trans} 的情况。此时两个直接子推导的结论
分别是 \(S\subtype U\) 和 \(U\subtype T\)。下面按照这两个子推导各自的
末条规则分类讨论；情形名称中斜线左侧表示 \(S\subtype U\) 推导的末条规则，
斜线右侧表示 \(U\subtype T\) 推导的末条规则。各情形的目标都是绕过中间类型
\(U\)，重新构造 \(S\subtype T\) 的无传递规则推导。

\emph{情形 任意规则／\textsc{S-Top}。}右侧推导以 \textsc{S-Top}
结束，所以 \(T=\tapltype{Top}\)。无论 \(S\) 是什么，直接使用
\textsc{S-Top} 就得到 \(S\subtype T\)。

\emph{情形 \textsc{S-Top}／任意规则。}左侧推导以 \textsc{S-Top}
结束，所以 \(U=\tapltype{Top}\)。先对右侧子推导使用归纳假设，设它已经
不含传递规则。检查其末条规则：自反规则已经消去，箭头规则和记录规则又都要求
左侧类型分别是箭头或记录，因此只有 \textsc{S-Top} 可能适用。于是
\(T=\tapltype{Top}\)，结论仍由 \textsc{S-Top} 得到。

\emph{情形 \textsc{S-Arrow}／\textsc{S-Arrow}。}写
\(S=S_1\to S_2\)、\(U=U_1\to U_2\)、\(T=T_1\to T_2\)。两个子推导
给出
\[
 U_1\subtype S_1,\quad S_2\subtype U_2,\quad
 T_1\subtype U_1,\quad U_2\subtype T_2
\]
为了直接应用 \textsc{S-Arrow}，需要得到 \(T_1\subtype S_1\) 和
\(S_2\subtype T_2\)。分别把上面两列中的推导用 \textsc{S-Trans} 接起来，
即可暂时构造出这两个判断的推导。它们都严格小于原来以传递规则结束的整个推导，
所以归纳假设可以再次消去其中的 \textsc{S-Trans}。最后应用
\textsc{S-Arrow}，便得到不含传递规则的
\(S_1\to S_2\subtype T_1\to T_2\) 推导。

\translatornote{这里的规模递减可以按规则节点数直接核对。原推导包含四份处理
参数和结果类型的内部推导，外面还有两次 \textsc{S-Arrow} 和最末端的一次
\textsc{S-Trans}，共多出三个规则节点。重新配对以后，两份新推导合起来仍使用
原来的四份内部推导，但只新加两次 \textsc{S-Trans}，总节点数已经少了一个；
分别来看，每份新推导只使用其中两份内部推导，当然也都严格小于原推导。归纳假设
正是分别用于这两份较小的新推导，所以临时加入传递规则并不会形成循环论证。}

\emph{情形 \textsc{S-Rcd}／\textsc{S-Rcd}。}证明相同：对目标记录中的
每个字段，先用 \textsc{S-Trans} 连接左右两个记录子推导给出的字段类型关系。
所得新推导都严格小于原推导，因此可以用归纳假设消去其中的传递规则；最后应用
\textsc{S-Rcd}，重新组成不含传递规则的记录子类型推导。

其余组合只有 \textsc{S-Arrow}／\textsc{S-Rcd} 和
\textsc{S-Rcd}／\textsc{S-Arrow}；它们会要求中间类型 \(U\) 同时是箭头
类型与记录类型，因此不可能发生。

\taplsolutionbacklink{lem:16.1.2}{返回引理16.1.2}
\end{taplauthorsolutionentrybox}

\begin{taplauthorsolutionentrybox}
\phantomsection
\label{sol:author:16.1.3}
\textbf{16.1.3 解答}

加入 \(\tapltype{Bool}\) 后，引理第一项应改为：“每个类型 \(S\) 的
\(S\subtype S\) 都能在不使用 \textsc{S-Refl} 的情况下导出，只有
\(S=\tapltype{Bool}\) 时例外。”另一种做法是在子类型定义中显式加入
\[
  \tapltype{Bool}\subtype\tapltype{Bool}
\]
再证明扩充后的系统可以删去 \textsc{S-Refl}。

\translatornote{\(\tapltype{Bool}\) 在这里成为例外，只因为它是没有内部结构的
基本类型。删去 \textsc{S-Refl} 后，\textsc{S-Top} 只能导出
\(\tapltype{Bool}\subtype\tapltype{Top}\)，\textsc{S-Arrow} 只适用于箭头
类型，\textsc{S-Rcd} 只适用于记录类型；因此没有规则能够导出
\(\tapltype{Bool}\subtype\tapltype{Bool}\)。每加入一种新的基本类型，都要为
它保留这样的自身子类型规则，或者继续使用通用的 \textsc{S-Refl}。}

\taplsolutionbacklink{ex:16.1.3}{返回练习16.1.3}
\end{taplauthorsolutionentrybox}

\begin{taplauthorsolutionentrybox}
\phantomsection
\label{sol:author:16.2.5}
\textbf{16.2.5 解答}

要证明的是：对任意声明式推导
\(\Gamma\vdash t:T\)，都能构造一份算法推导
\(\Gamma\algturnstile t:S\)，且 \(S\subtype T\)。对给定的声明式推导作归纳，
按其最后使用的规则分为以下情形。在每个情形中，归纳假设均应用于
末条规则的直接子推导。

\textbf{情形 \textsc{T-Var}。}已知 \(x:T\in\Gamma\)。由
\textsc{TA-Var} 得到 \(\Gamma\algturnstile x:T\)。取 \(S=T\)；由子类型关系的
自反性，\(S\subtype T\)。

\textbf{情形 \textsc{T-Abs}。}声明式推导的直接子推导为
\[
  \Gamma,x:T_1\vdash t_2:T_2
\]
结论为
\(\Gamma\vdash\lambda x:T_1.\,t_2:T_1\to T_2\)。归纳假设给出某个
\(S_2\subtype T_2\)，使
\(\Gamma,x:T_1\algturnstile t_2:S_2\)。使用 \textsc{TA-Abs} 得到
\[
  \Gamma\algturnstile\lambda x:T_1.\,t_2:T_1\to S_2
\]
又因为 \(T_1\subtype T_1\) 且 \(S_2\subtype T_2\)，规则
\textsc{S-Arrow} 给出
\[
  T_1\to S_2\subtype T_1\to T_2
\]
所以算法规则为这个抽象算出了声明式类型的一个子类型。

\textbf{情形 \textsc{T-App}。}声明式推导的两个直接子推导为
\[
  \Gamma\vdash t_1:T_{11}\to T_{12}
  \qquad
  \Gamma\vdash t_2:T_{11}
\]
结论为 \(\Gamma\vdash t_1\,t_2:T_{12}\)。分别应用归纳假设，得到
\[
\begin{gathered}
  \Gamma\algturnstile t_1:S_1,
  \qquad S_1\subtype T_{11}\to T_{12},\\
  \Gamma\algturnstile t_2:S_2,
  \qquad S_2\subtype T_{11}
\end{gathered}
\]
由子类型反演，\(S_1\) 必须也是箭头类型。因此可以写成
\(S_1=S_{11}\to S_{12}\)，并且有
\[
  T_{11}\subtype S_{11}
  \qquad\text{和}\qquad
  S_{12}\subtype T_{12}
\]
实参的算法类型满足 \(S_2\subtype T_{11}\)，而参数类型又满足
\(T_{11}\subtype S_{11}\)，所以传递性给出
\[
  S_2\subtype S_{11}
\]
再由算法子类型关系的完备性，可将该判断改写为
\(\algturnstile S_2\subtype S_{11}\)。这正是 \textsc{TA-App} 所需的子类型前提，
因而可以导出
\[
  \Gamma\algturnstile t_1\,t_2:S_{12}
\]
前面已经得到 \(S_{12}\subtype T_{12}\)，因此应用情形完成。

\textbf{情形 \textsc{T-Rcd}。}对每个字段的直接子推导
\(\Gamma\vdash t_i:T_i\) 应用归纳假设，得到某个
\(S_i\subtype T_i\)，使 \(\Gamma\algturnstile t_i:S_i\)。使用
\textsc{TA-Rcd} 可得
\[
  \Gamma\algturnstile
  \{l_i=t_i\}^{i\in1..n}:\{l_i:S_i\}^{i\in1..n}
\]
由记录子类型规则，每个 \(S_i\subtype T_i\) 共同给出
\[
  \{l_i:S_i\}^{i\in1..n}\subtype\{l_i:T_i\}^{i\in1..n}
\]
因此记录情形成立。

\textbf{情形 \textsc{T-Proj}。}设声明式前提为
\[
  \Gamma\vdash t_1:\{l_i:T_i\}^{i\in1..n}
\]
而结论是 \(\Gamma\vdash t_1.l_j:T_j\)。归纳假设给出
\(\Gamma\algturnstile t_1:S\)，且
\[
  S\subtype\{l_i:T_i\}^{i\in1..n}
\]
由子类型反演，\(S\) 必须是一个记录类型，并且至少包含字段
\(l_j:S_j\)，其中 \(S_j\subtype T_j\)。因此 \textsc{TA-Proj} 导出
\[
  \Gamma\algturnstile t_1.l_j:S_j
\]
而 \(S_j\subtype T_j\)，正好得到所需结论。

\textbf{情形 \textsc{T-Sub}。}设末条规则的直接子推导为
\(\Gamma\vdash t:U\)，并且 \(U\subtype T\)。归纳假设给出某个
\(S\subtype U\)，使 \(\Gamma\algturnstile t:S\)。由子类型关系的传递性，
\(S\subtype T\)。同一份算法推导因此就满足定理结论。

\taplsolutionbacklink{thm:16.2.5}{返回定理16.2.5}
\end{taplauthorsolutionentrybox}

\begin{taplauthorsolutionentrybox}
\phantomsection
\label{sol:author:16.2.6}
\textbf{16.2.6 解答}

若删去 \textsc{S-Arrow}，项
\(\lambda x:\{a:\tapltype{Nat}\}.\,x\) 在声明式规则下同时具有
\[
\{a:\tapltype{Nat}\}\to\{a:\tapltype{Nat}\}
\quad\text{和}\quad
\{a:\tapltype{Nat}\}\to\tapltype{Top}
\]
两种类型。删去箭头子类型规则后，这两个类型不可比较，也不存在同时位于二者
之下的类型，所以该项没有最小类型。

\taplsolutionbacklink{ex:16.2.6}{返回练习16.2.6}
\end{taplauthorsolutionentrybox}

\begin{taplauthorsolutionentrybox}
\phantomsection
\label{sol:author:16.3.2}
\textbf{16.3.2 解答}

以下两个算法相互递归。并算法总会返回一个类型；交算法在不存在交时失败。
\[
S\sqcup T=
\left\{
\begin{array}{@{}ll@{}}
\tapltype{Bool}
  & \text{若 }S=T=\tapltype{Bool}\\[2pt]
M_1\to J_2
  & \begin{aligned}
      &\text{若 }S=S_1\to S_2,\quad T=T_1\to T_2,\\
      &S_1\sqcap T_1=M_1,\quad S_2\sqcup T_2=J_2
    \end{aligned}\\[5pt]
\{j_i:J_i\}^{i\in1..q}
  & \begin{aligned}
      &\text{若 }S=\{k_j:S_j\}^{j\in1..m},\quad
        T=\{l_i:T_i\}^{i\in1..n},\\
      &\{j_i\}^{i\in1..q}
        =\{k_j\}^{j\in1..m}\cap\{l_i\}^{i\in1..n},\\
      &S_j\sqcup T_i=J_i
        \quad\text{对每个 }j_i=k_j=l_i
    \end{aligned}\\[7pt]
\tapltype{Top}
  & \text{其他情形}
\end{array}
\right.
\]
\[
S\sqcap T=
\left\{
\begin{array}{@{}ll@{}}
S
  & \text{若 }T=\tapltype{Top}\\[2pt]
T
  & \text{若 }S=\tapltype{Top}\\[2pt]
\tapltype{Bool}
  & \text{若 }S=T=\tapltype{Bool}\\[2pt]
J_1\to M_2
  & \begin{aligned}
      &\text{若 }S=S_1\to S_2,\quad T=T_1\to T_2,\\
      &S_1\sqcup T_1=J_1,\quad S_2\sqcap T_2=M_2
    \end{aligned}\\[5pt]
\{m_i:M_i\}^{i\in1..q}
  & \begin{aligned}
      &\text{若 }S=\{k_j:S_j\}^{j\in1..m},\quad
        T=\{l_i:T_i\}^{i\in1..n},\\
      &\{m_i\}^{i\in1..q}
        =\{k_j\}^{j\in1..m}\cup\{l_i\}^{i\in1..n},\\
      &S_j\sqcap T_i=M_i
        \quad\text{对每个 }m_i=k_j=l_i,\\
      &M_i=S_j
        \quad\text{若 }m_i=k_j\text{ 只出现在 }S\text{ 中},\\
      &M_i=T_i
        \quad\text{若 }m_i=l_i\text{ 只出现在 }T\text{ 中}
    \end{aligned}\\[9pt]
\mathsf{fail}
  & \text{其他情形}
\end{array}
\right.
\]

\translatornote{箭头类型的并在参数位置使用交，正是参数逆变的结果。设
\(U_1\to U_2\) 是 \(S_1\to S_2\) 与 \(T_1\to T_2\) 的公共超类型。
由 \textsc{S-Arrow}，前一个子类型判断要求
\(U_1\subtype S_1\) 且 \(S_2\subtype U_2\)，后一个要求
\(U_1\subtype T_1\) 且 \(T_2\subtype U_2\)。因此，公共超类型的参数类型
\(U_1\) 必须是 \(S_1,T_1\) 的共同下界，结果类型 \(U_2\) 则必须是
\(S_2,T_2\) 的共同上界。为了得到最小的公共函数超类型，前者应取最大
公共下界，即交；后者应取最小公共上界，即并。这就得到
\[
  (S_1\to S_2)\sqcup(T_1\to T_2)
  =(S_1\sqcap T_1)\to(S_2\sqcup T_2)
\]
例如，若某个函数可能具有
\(\{x:\tapltype{Bool}\}\to\tapltype{Bool}\) 或
\(\{y:\tapltype{Bool}\}\to\tapltype{Bool}\) 中的任一类型，那么只有同时含有
\(x\) 和 \(y\) 字段的实参才肯定能被两种函数接受。因此，这两个函数类型的并是
\[
  \{x:\tapltype{Bool},y:\tapltype{Bool}\}\to\tapltype{Bool}
\]
求交时情形对偶：参数位置取并，结果位置取交。}

记录交中，公共字段取字段类型的交；只出现在一侧的字段原样保留。任何递归
调用失败时，整次交计算失败。递归调用时两种输入类型的总大小严格减小，因此
算法不会发散。

并算法的箭头分支会调用交算法计算参数类型的交；若这次调用失败，并算法就继续
落入最后的默认分支，返回 \(\tapltype{Top}\)。例如，
\[
 (\tapltype{Bool}\to\tapltype{Top})\sqcup
 (\{\}\to\tapltype{Top})=\tapltype{Top}
\]
两种算法都不会发散，因为每次递归调用都会减小输入 \(S,T\) 的总大小。交算法
虽然可能失败，但只要两个输入存在共同下界，它就一定成功。

\setcounter{taplappendixlemma}{9}
\begin{taplappendixlemma}
\label{lem:appendix-a10}
若 \(L\subtype S\) 且 \(L\subtype T\)，则存在某个 \(M\)，使
\(S\sqcap T=M\)。
\end{taplappendixlemma}

\begin{proof}
对 \(S\) 的大小归纳（等价地，也可对 \(T\) 的大小归纳），并按 \(S,T\)
的外层形式分类。若其中一个是 \(\tapltype{Top}\)，交算法的前两个分支之一
适用，结果分别为另一个输入。若 \(S,T\) 的外层形式不同，则子类型反演引理
会对共同下界 \(L\) 的形式提出互不相容的要求。例如，\(S\) 是箭头类型时，
\(L\) 也必须是箭头类型；\(T\) 是记录类型时，\(L\) 又必须是记录类型。\footnote{
严格地说，\cref{lem:15.3.2}没有讨论 \(\tapltype{Bool}\)。补上的布尔情形只需
说明：\(\tapltype{Bool}\) 的子类型只有它自身。}因此只剩三种情形。

若 \(S=T=\tapltype{Bool}\)，交算法的第三个分支立即返回
\(\tapltype{Bool}\)。

若 \(S=S_1\to S_2\) 且 \(T=T_1\to T_2\)，并算法的全定义性保证第一次
递归调用 \(S_1\sqcup T_1\) 返回某个 \(J_1\)。子类型反演又说明
\(L=L_1\to L_2\)，并且 \(L_2\subtype S_2\)、\(L_2\subtype T_2\)。
于是 \(L_2\) 是 \(S_2,T_2\) 的共同下界，归纳假设保证
\(S_2\sqcap T_2\) 成功并返回某个 \(M_2\)。所以
\(S\sqcap T=J_1\to M_2\)。

最后，设 \(S=\{k_j:S_j\}^{j\in1..m}\)，
\(T=\{l_i:T_i\}^{i\in1..n}\)。由子类型反演，\(L\) 必是记录类型，而且
包含 \(S,T\) 中出现的全部标签。对同时出现在 \(S,T\) 中的每个标签，
\(L\) 中对应字段类型都是两侧字段类型的共同子类型；归纳假设因此保证交算法
对所有公共字段的递归调用都成功。只出现在一侧的字段无需递归求交，故整次记录
求交成功。
\end{proof}

下面验证这些定义确实计算出并与交。论证分成两部分：
\taplnamedref{prop:appendix-a11}{命题}说明求得的交是 \(S,T\) 的下界、求得的
并是它们的上界；\taplnamedref{prop:appendix-a12}{命题}说明求得的交大于每个
共同下界，求得的并小于每个共同上界。

\taplstatementliststart
\begin{taplappendixproposition}
\label{prop:appendix-a11}
\taplstatementlistbreak
\begin{enumerate}
\item 若 \(S\sqcup T=J\)，则 \(S\subtype J\) 且 \(T\subtype J\)。
\item 若 \(S\sqcap T=M\)，则 \(M\subtype S\) 且 \(M\subtype T\)。
\end{enumerate}
\end{taplappendixproposition}

\begin{proof}
对 \(S\sqcup T=J\) 或 \(S\sqcap T=M\) 的“推导”大小作直接归纳；这里的
大小就是为了算出 \(J\) 或 \(M\) 而对两个定义进行的递归调用次数。
\end{proof}

\taplstatementliststart
\begin{taplappendixproposition}
\label{prop:appendix-a12}
\taplstatementlistbreak
\begin{enumerate}
\item 设 \(S\sqcup T=J\)，且对某个 \(U\) 有
  \(S\subtype U\) 与 \(T\subtype U\)。则 \(J\subtype U\)。
\item 设 \(S\sqcap T=M\)，且对某个 \(L\) 有
  \(L\subtype S\) 与 \(L\subtype T\)。则 \(L\subtype M\)。
\end{enumerate}
\end{taplappendixproposition}

\begin{proof}
两部分同时对 \(S,T\) 的大小归纳；给定 \(S,T\) 后，依次证明两部分。

先证第一项，并按 \(U\) 的形式分类。若 \(U=\tapltype{Top}\)，则任何 \(J\)
都满足 \(J\subtype U\)。若 \(U=\tapltype{Bool}\)，子类型反演说明
\(S=T=\tapltype{Bool}\)，所以 \(J=\tapltype{Bool}\)。

若 \(U=U_1\to U_2\)，子类型反演给出
\(S=S_1\to S_2\)、\(T=T_1\to T_2\)，并且
\[
 U_1\subtype S_1,\quad U_1\subtype T_1,\quad
 S_2\subtype U_2,\quad T_2\subtype U_2
\]
由归纳假设，\(S_1,T_1\) 的交 \(M_1\) 位于 \(U_1\) 之上，即
\(U_1\subtype M_1\)；\(S_2,T_2\) 的并 \(J_2\) 位于 \(U_2\) 之下，
即 \(J_2\subtype U_2\)。应用 \textsc{S-Arrow}，得到
\(M_1\to J_2\subtype U_1\to U_2\)。

若 \(U\) 是记录类型，子类型反演说明 \(S,T\) 也都是记录类型；\(S,T\)
的标签集合都包含 \(U\) 的标签集合，而且 \(U\) 每个字段的类型都是两侧对应
字段类型的超类型。因此，\(S,T\) 的并至少包含 \(U\) 的全部标签；对每个
公共字段，归纳假设说明并中的字段类型是 \(U\) 对应字段类型的子类型。
由 \textsc{S-Rcd} 得 \(J\subtype U\)。

再证第二项，并按 \(S,T\) 的形式分类。若其中一个是 \(\tapltype{Top}\)，
交就是另一个输入，结论立即成立。两个非最大类型具有不同外层形式的情形，由
\taplnamedref{lem:appendix-a10}{引理}的论证可知不可能发生。若二者都是
\(\tapltype{Bool}\)，结论同样
立即成立。

若 \(S=S_1\to S_2\)、\(T=T_1\to T_2\)，子类型反演给出
\(L=L_1\to L_2\)，并且
\[
 S_1\subtype L_1,\quad T_1\subtype L_1,\quad
 L_2\subtype S_2,\quad L_2\subtype T_2
\]
由归纳假设，\(S_1,T_1\) 的并 \(J_1\) 满足 \(J_1\subtype L_1\)，
\(S_2,T_2\) 的交 \(M_2\) 满足 \(L_2\subtype M_2\)。再由
\textsc{S-Arrow} 得 \(L_1\to L_2\subtype J_1\to M_2\)。

若 \(S,T\) 都是记录类型，子类型反演说明 \(L\) 也是记录类型，并包含
\(S,T\) 的全部标签（还可能更多），对应字段类型满足所需子类型关系。对交
\(M\) 中的每个标签 \(m_l\)，分三种情况。若它同时出现在 \(S,T\) 中，
其在 \(M\) 中的类型是两侧字段类型的交，归纳假设说明 \(L\) 中对应字段类型
是它的子类型。若 \(m_l\) 只出现在 \(S\) 中，则它在 \(M\) 中沿用 \(S\)
中的类型，而我们已知 \(L\) 中对应字段类型更小。只出现在 \(T\) 中的情况
相同。于是由 \textsc{S-Rcd} 得 \(L\subtype M\)。
\end{proof}

\taplsolutionbacklink{prop:16.3.2}{返回命题16.3.2}
\end{taplauthorsolutionentrybox}

\begin{taplauthorsolutionentrybox}
\phantomsection
\label{sol:author:16.3.3}
\textbf{16.3.3 解答}

最小类型是 \(\tapltype{Top}\)，即 \(\tapltype{Bool}\) 与 \(\{\}\)
的并。不过，这个项能够通过类型检查可视为语言的一个弱点：
\(\tapltype{Top}\) 型值没有可供使用的操作，程序员通常不会有意写出这样的
条件式。一种做法是从系统中删除 \(\tapltype{Top}\)，并把求并改成偏操作；
另一种做法是保留 \(\tapltype{Top}\)，但让类型检查器在项的最小类型为
\(\tapltype{Top}\) 时发出警告。

\taplsolutionbacklink{ex:16.3.3}{返回练习16.3.3}
\end{taplauthorsolutionentrybox}

\begin{taplauthorsolutionentrybox}
\phantomsection
\label{sol:author:16.3.4}
\textbf{16.3.4 解答}

处理不变的 \(\tapltype{Ref}\) 很直接，只需在并、交算法中分别加入一条分支：
\[
\begin{array}{rcl}
\tapltype{Ref}\ S_1\sqcup\tapltype{Ref}\ T_1
 &=&\tapltype{Ref}\ T_1
 \quad\text{若 }S_1\subtype T_1\text{ 且 }T_1\subtype S_1\\[2pt]
\tapltype{Ref}\ S_1\sqcap\tapltype{Ref}\ T_1
 &=&\tapltype{Ref}\ T_1
 \quad\text{若 }S_1\subtype T_1\text{ 且 }T_1\subtype S_1
\end{array}
\]

\translatornote{这里沿用\taplsecref{15.5}{sec:15.5}中 \(\tapltype{Ref}\) 的不变子类型
规则。一个 \(\tapltype{Ref}\ S_1\) 既能读取又能写入：要把它当作
\(\tapltype{Ref}\ T_1\) 使用，读取要求 \(S_1\subtype T_1\)，写入则要求
\(T_1\subtype S_1\)。因此，只有在两个元素类型互为子类型时，
\(\tapltype{Ref}\ S_1\) 与 \(\tapltype{Ref}\ T_1\) 才能相互替代；此时任取其中一个，
都可作为二者的并或交。若两者不等价，直接取
\(\tapltype{Ref}(S_1\sqcup T_1)\) 会扩大允许写入的值的范围，不再类型安全。
此时求并会落入默认分支得到 \(\tapltype{Top}\)，求交则可能失败。}

把 \(\tapltype{Ref}\) 分解为 \(\tapltype{Source}\) 和
\(\tapltype{Sink}\) 后会遇到严重困难：子类型关系不再总有并或交。例如，
\(\tapltype{Ref}\{a:\tapltype{Nat},b:\tapltype{Bool}\}\) 与
\(\tapltype{Ref}\{a:\tapltype{Nat}\}\) 同时是
\(\tapltype{Source}\{a:\tapltype{Nat}\}\) 和
\(\tapltype{Sink}\{a:\tapltype{Nat},b:\tapltype{Bool}\}\) 的子类型，
但后两个类型没有共同下界。

最简单的解决方式之一，是只加入 \(\tapltype{Source}\) 或只加入
\(\tapltype{Sink}\)。许多应用只需其中一个。例如，
\taplsecref{18.12}{sec:18.12}中的类实现只需 \(\tapltype{Source}\)；并发
语言中的服务器进程则可能只需要向外传递发送能力，即 \(\tapltype{Sink}\)；
接收能力只由服务器自身持有。

若系统只加入 \(\tapltype{Source}\)，并算法在加入上面的 \(\tapltype{Ref}\)
分支后，再加入以下四种组合，仍然是完备的：
\[
\begin{array}{rcl}
\tapltype{Ref}\ S_1\sqcup\tapltype{Ref}\ T_1
 &=&\tapltype{Source}\ J_1
 \quad\text{若 }S_1\sqcup T_1=J_1\\
\tapltype{Source}\ S_1\sqcup\tapltype{Source}\ T_1
 &=&\tapltype{Source}\ J_1
 \quad\text{若 }S_1\sqcup T_1=J_1\\
\tapltype{Ref}\ S_1\sqcup\tapltype{Source}\ T_1
 &=&\tapltype{Source}\ J_1
 \quad\text{若 }S_1\sqcup T_1=J_1\\
\tapltype{Source}\ S_1\sqcup\tapltype{Ref}\ T_1
 &=&\tapltype{Source}\ J_1
 \quad\text{若 }S_1\sqcup T_1=J_1
\end{array}
\]
交算法也加入类似分支。

另一种方案由 Hennessy 与 Riely（1998）提出：让引用类型接受两个参数。
\(\tapltype{Ref}\ S\ T\) 允许写入 \(S\) 型值、
读出 \(T\) 型值；它对第一个参数逆变，对第二个参数协变。此时可把
\(\tapltype{Sink}\ S\) 定义为
\(\tapltype{Ref}\ S\ \tapltype{Top}\)，把
\(\tapltype{Source}\ T\) 定义为
\(\tapltype{Ref}\ \tapltype{Bot}\ T\)。

\taplsolutionbacklink{ex:16.3.4}{返回练习16.3.4}
\end{taplauthorsolutionentrybox}

\begin{taplauthorsolutionentrybox}
\phantomsection
\label{sol:author:16.4.1}
\textbf{16.4.1 解答}

需要。合适的规则是
\[
\inferrule*[right=\textsc{TA-If}]
 {\Gamma\algturnstile t_1:T_1\\T_1=\tapltype{Bot}\\
  \Gamma\algturnstile t_2:T_2\\
  \Gamma\algturnstile t_3:T_3\\T_2\sqcup T_3=T}
 {\Gamma\algturnstile
  \mathsf{if}\ t_1\ \mathsf{then}\ t_2\ \mathsf{else}\ t_3:T}
\]
另一条看似自然的候选规则会直接把整个条件式赋为 \(\tapltype{Bot}\)：
\[
\inferrule*[right=\textsc{TA-If}]
 {\Gamma\algturnstile t_1:T_1\\T_1=\tapltype{Bot}\\
  \Gamma\algturnstile t_2:T_2\\
  \Gamma\algturnstile t_3:T_3}
 {\Gamma\algturnstile
  \mathsf{if}\ t_1\ \mathsf{then}\ t_2\ \mathsf{else}\ t_3:
  \tapltype{Bot}}
\]
这条规则从运行行为看是安全的：\(\tapltype{Bot}\) 没有值，条件项不可能
求得 \(\tapltrue\) 或 \(\taplfalse\)，因而整个条件式不会产生普通结果。但它可能给
某些项赋予声明式规则无法赋予的类型，从而破坏算法类型关系的可靠性定理
\cref{thm:16.2.4}。

\translatornote{例如，设上下文中有 \(x:\tapltype{Bot}\)，考虑
\(\mathsf{if}\ x\ \mathsf{then}\ \tapltrue\ \mathsf{else}\ \taplfalse\)。上面的候选规则会把
整个条件式赋为 \(\tapltype{Bot}\)。声明式系统可以先由
\(\tapltype{Bot}\subtype\tapltype{Bool}\) 把 \(x\) 当作布尔值，再根据两个分支得到整个
条件式的类型 \(\tapltype{Bool}\)；它无法反向把该类型降为 \(\tapltype{Bot}\)。
因此，这个项并非完全无法由声明式系统赋型；问题在于，候选算法规则赋予它的
\(\tapltype{Bot}\) 无法由声明式规则推导出来。}

\taplsolutionbacklink{ex:16.4.1}{返回练习16.4.1}
\end{taplauthorsolutionentrybox}
